\label{chap:introduction}
Students in higher education face growing challenges in planning their studies. Curricula involve complex prerequisites, electives, and institutional rules that students must reconcile. While institutions publish exemplary study plans, these assume an idealised progression; a single failed examination or postponed module immediately necessitates an individualised schedule~\cite{judel_supporting_2023}. Varying study speeds and limited access to guidance compound this issue. Advisors struggle to provide evidence-based scheduling advice at scale~\cite{schulte_large_2017}, which is critical given that the timing and sequencing of courses are associated with time to degree, and less strongly with final grade~\cite{morsy_study_2019}.

Existing approaches address these issues but remain limited. Graph-based systems focus mainly on visualising curriculum structure and the \gterm{dependency}{dependencies} it encodes~\cite{auvinen_stops_2014, rollande_graph_2013, trippel_developing_2025}, while table-based planning tools emphasise semester-level scheduling~\cite{hirmer_requirements_2022, judel_supporting_2023, srisamutr_course_2018}; where both appear in one tool, the dependency view stays local to a single module. This motivates investigating a hybrid approach: integrating a graph and a table into a single planning workflow so students can understand structure and dependencies (graph) and manage their schedule and workload (table).
Two further strands turn student data into feedback. Educational Dashboards (ED) report a student's standing~\cite{schwendimann2017, verbert_learning_2020, loboda_mastery_2014}, while Educational Recommender Systems (ERS) suggest which courses to take next~\cite{bhumichitr_recommender_2017, esteban_helping_2020, al-badarenah_automated_2016}. Both are well established, and both are almost always built separately~\cite{bodily_review_2017}. How the two should be combined within a student's own planning workflow therefore remains open. This thesis puts that question to students, in the formative study of Chapter~\ref{chap:formative-study}, and derives the answer from what they report.

\section{Research Questions}
\label{sec:research-questions}

Building on the challenges and gaps outlined above, this thesis investigates the design, implementation, and use of a hybrid graph-table study planner, built for and evaluated against two degree programmes at the Faculty of Informatics of the Vienna University of Technology (TU~Wien): the Bachelor programme in Computer Science\footnote{\url{https://informatics.tuwien.ac.at/bachelor/curriculum-ue033521.pdf}} and the Master programme in Software Engineering\footnote{\url{https://informatics.tuwien.ac.at/master/curriculum-ue066937.pdf}}. We formulate three research questions:

\textbf{Research Question 1} \emph{What are students' needs and resulting requirements for a hybrid graph-table interface that supports programme-level study planning and understanding of curriculum structure and course dependencies?}

This question is addressed by the formative study, conducted with an interactive but non-functional prototype (Chapter~\ref{chap:formative-study}), whose findings are translated into a ranked features-and-requirements specification in the same chapter (Section~\ref{chap:from-themes-to-requirements}).
Building on these requirements, the second sub-question concerns the design and construction of the tool itself:

\textbf{Research Question 2} \emph{How can we design and implement a hybrid interface that combines graph-based visualisation with table-based planning to help students plan their studies and understand curriculum structure and course dependencies, based on those requirements?}

This question is addressed by the Design chapter (Chapter~\ref{chap:design}), which derives each design decision from a requirement, prior evidence, or an explicitly stated assumption, and by the Implementation chapter (Chapter~\ref{chap:implementation}), which realises that design as a working system.
Finally, the third question asks how the implemented system is used by the students it was built for, and what its graph contributes to that use:

\textbf{Research Question 3} \emph{How do students plan with the hybrid interface, and what does the graph-based view contribute to their planning?}

The evaluation reads this question through three lenses: the planning process as the recordings show it, the plans students produce, measured against their brief and against what students and the tool report as complete, and the satisfaction students report. This question is addressed by an evaluation study with students (Chapter~\ref{chap:evaluation}), interpreted in the Discussion (Chapter~\ref{chap:discussion}) and answered in the Conclusion (Chapter~\ref{chap:conclusion}).

\section{Contributions}
\label{sec:contributions}

This thesis makes the following contributions:

\TDrev{Re-check the contributions once Evaluation and Discussion are final}{Kept open at your request. Contribution one now states the specification's size, and contributions three and four track Chapter~\ref{chap:evaluation}'s two headline results. One thing to settle when Chapters~\ref{chap:evaluation} and~\ref{chap:discussion} are closed: contribution two says the tool \enquote{embeds explainable recommendations}, which is true of what was designed, built and evaluated: all six families were live during the sessions (settled 2026-09-05; Chapter~\ref{chap:implementation}, Section~\ref{sec:impl-recommender}). What the study cannot speak to is the quality of the suggestions, which OOS5 already places out of scope, so the contribution needs no narrowing on that account.}
\begin{itemize}
    \item \textbf{A requirements specification for study planning.} A ranked specification of 13 features carrying 55 numbered requirements for a study-planning tool, each traceable through a theme to the formative interview evidence that motivated it (Section~\ref{chap:from-themes-to-requirements}).
    \item \textbf{A hybrid graph-table planning tool.} The design and implementation of a web-based tool that unites a curriculum graph and a schedule-managing table over a single shared plan, checks that plan against the regulations of two degree programmes in real time, and embeds explainable recommendations within the planning workflow (Chapters~\ref{chap:design}--\ref{chap:implementation}).
    \item \textbf{An account of what a curriculum graph contributes to planning.} Evidence, from a within-subjects study of eleven students observed at the level of the interface itself, that the graph enters the planning process at one point: it replaces the Course Catalogue as the surface on which a named gap is turned into candidate courses, and it leaves untouched the step at which a course is committed to a semester, because it renders no semester dimension against which that commitment could be judged (Chapter~\ref{chap:evaluation}).
    \item \textbf{Evidence that a completion signal students trust can certify a plan it has not checked.} Plans that missed their brief were reported as successes by the students who produced them, and the tool agreed, because the constraints it does not model are silent rather than wrong. This bears on any planning tool that guides through a checklist (Chapter~\ref{chap:evaluation}).
\end{itemize}

The last two contributions rest on two records. The account of how students planned is read from a frame-sampled log of the interface state; the catalogue of usability problems and positive findings is coded from that log and from the session transcripts together, each code carrying a participant-level frequency and mapped to an ISO~9241-11 dimension~\cite{iso9241_11}; the questionnaires form a third, descriptive layer, which locates the weaknesses without carrying a finding of its own. Because the evaluation used a fixed scenario order and eleven analysed participants, this evidence is descriptive rather than causal: it describes how students approached planning, but does not establish that a particular interface component caused a particular behaviour. Which of these findings are specific to the graph-table hybrid and which might carry to other planning tools is examined in Chapter~\ref{chap:discussion} (Section~\ref{sec:disc-generalisability}).

\section{Outlook}
\label{sec:outlook}

The remainder of this thesis is organised as follows. Chapter~\ref{chap:methodology} sets out the research paradigm, process, and scope that frame the work. Chapter~\ref{chap:related-work} reviews related systems and establishes the research gap. Chapter~\ref{chap:formative-study} reports the formative study and translates its findings into a features-and-requirements specification, answering Research Question~1. Chapter~\ref{chap:design} and Chapter~\ref{chap:implementation} design and build the hybrid tool, answering Research Question~2. Chapter~\ref{chap:evaluation} evaluates the tool with students, answering Research Question~3, and Chapter~\ref{chap:discussion} interprets what it found, relates it to prior work and sets out the limitations that bound it. Chapter~\ref{chap:conclusion} states the answer to each of the three research questions in turn (Section~\ref{sec:disc-rq-answers}), restates the contributions, and names the work that follows.
